\documentclass[../main.tex]{subfiles}
\begin{document}

\chapter{Arquitectura del laboratorio}
\label{C:arquitectura}

\section{Entorno de experimentación}

\subsection{Hardware del host}

El laboratorio completo se ejecuta sobre un equipo de sobremesa personal con las
siguientes especificaciones:

\begin{table}[ht]
  \centering
  \caption{Especificaciones del equipo \textit{host} de experimentación.}
  \label{tab:hardware}
  \rowcolors{2}{tabla-par}{white}
  \begin{tabularx}{\textwidth}{X X}
    \toprule
    \rowcolor{tabla-cabecera}
    \textcolor{white}{\textbf{Componente}} &
    \textcolor{white}{\textbf{Especificación}} \\
    \midrule
    Procesador        & AMD Ryzen (arquitectura x86\_64) \\
    Memoria RAM       & 16\,GB DDR4 \\
    Almacenamiento    & SSD NVMe \\
    Sistema operativo & CachyOS (basado en Arch Linux) \\
    Kernel            & linux-cachyos 7.1.3-2 \\
    Docker Engine     & 27.x \\
    Docker Compose    & v2.x \\
    \bottomrule
  \end{tabularx}
\end{table}

El consumo de memoria en reposo del \emph{stack} completo es de aproximadamente
6\,GB, distribuidos principalmente entre el \emph{indexer} OpenSearch (2\,GB de
\emph{heap} fijo) y el \emph{manager} Wazuh (512\,MB).

\subsection{Sistema operativo: CachyOS}

CachyOS es una distribución Linux basada en Arch Linux orientada al rendimiento,
que incluye el \emph{kernel} personalizado linux-cachyos con parches para mejorar
la latencia en cargas de trabajo interactivas. Entre estos parches se encuentran
el planificador \gls{bore} y modificaciones en la gestión de memoria y
caché de \gls{cpu}.

Estas optimizaciones tienen una consecuencia relevante para el laboratorio: el
\emph{kernel} linux-cachyos bloquea el acceso al subsistema de auditoría del
\emph{kernel} (\emph{netlink audit}) desde \emph{namespaces} de contenedor, incluso
con \texttt{privileged:~true}. Esta limitación y la solución adoptada se documentan
en la sección~\ref{sec:limitaciones}.

\section{Visión general del diseño}

El laboratorio se implementa como un conjunto de contenedores Docker orquestados
mediante Docker Compose. La decisión de usar contenedores en lugar de máquinas
virtuales responde a tres criterios:

\begin{enumerate}
  \item \textbf{Reproducibilidad}: los contenedores garantizan un entorno idéntico
  en cualquier máquina con Docker instalado. Cualquier investigador puede reproducir
  el experimento clonando el repositorio y ejecutando \texttt{docker compose up}.

  \item \textbf{Eficiencia de recursos}: el \emph{overhead} de los contenedores
  es significativamente menor que el de las \gls{vm}. El laboratorio opera con
  6\,GB de \gls{ram}, frente a los 16--32\,GB de una implementación equivalente
  con máquinas virtuales.

  \item \textbf{Facilidad de automatización}: Docker Compose define el estado
  completo de la infraestructura en un único fichero declarativo, facilitando la
  integración con el \emph{script} de orquestación y la reconfiguración entre
  ejecuciones.
\end{enumerate}

La infraestructura se organiza en una red privada \texttt{172.20.0.0/24} con IPs
estáticas asignadas a cada contenedor. La tabla~\ref{tab:contenedores} detalla la
función de cada contenedor.

\begin{table}[ht]
  \centering
  \caption{Contenedores del laboratorio y su dirección IP en la red \texttt{lab-net}.}
  \label{tab:contenedores}
  \rowcolors{2}{tabla-par}{white}
  \begin{tabularx}{\textwidth}{C{2.5cm} C{3cm} X}
    \toprule
    \rowcolor{tabla-cabecera}
    \textcolor{white}{\textbf{Contenedor}} &
    \textcolor{white}{\textbf{IP}} &
    \textcolor{white}{\textbf{Función}} \\
    \midrule
    wazuh-indexer   & 172.20.0.21 & OpenSearch: almacenamiento de alertas \\
    wazuh-manager   & 172.20.0.20 & Motor de correlación y reglas \\
    wazuh-dashboard & 172.20.0.22 & Interfaz web de visualización \\
    fileserver      & 172.20.0.10 & Víctima principal (servidor de ficheros) \\
    workstation     & 172.20.0.11 & Víctima secundaria (objetivo de E3) \\
    attacker        & 172.20.0.30 & Atacante basado en Kali Linux \\
    restic-sftp     & 172.20.0.41 & Destino de backups estrategia B \\
    minio           & 172.20.0.40 & Almacenamiento de objetos S3-compatible \\
    minio-init      & ---         & Inicialización de buckets (one-shot) \\
    \bottomrule
  \end{tabularx}
\end{table}

\section{Descripción de los contenedores}

\subsection{Infraestructura SIEM}

El \textbf{wazuh-indexer} ejecuta OpenSearch~4.7.5. Requiere
\texttt{vm.max\_map\_count}~$\geq$~262.144 en el \emph{kernel} del \emph{host}
y se configura con \emph{heap} fijo de 1\,GB (\texttt{-Xms1g} \texttt{-Xmx1g}) para
evitar pausas de recolección de basura de la \gls{jvm}.

El \textbf{wazuh-manager} es el núcleo del \gls{siem}: recibe los eventos de los
agentes, aplica las reglas de correlación y envía las alertas al \emph{indexer}
mediante Filebeat con \gls{tls} mutuo. La configuración activa \texttt{logall=yes}
y \texttt{logall\_json=yes} para archivar todos los eventos, no solo los que generan
alerta. Esto resultó fundamental durante la calibración del laboratorio para
verificar que los eventos llegaban correctamente aunque no disparasen ninguna regla.

El \textbf{wazuh-dashboard} proporciona la interfaz web de visualización y
administración, accesible en \url{https://localhost}.

\subsection{Sistemas víctima}

El \textbf{fileserver} es el objetivo principal de los ataques. Ejecuta Ubuntu~24.04
con los siguientes componentes:

\begin{itemize}
  \item Servidor \gls{ssh} con autenticación exclusivamente por clave pública
  (\texttt{PasswordAuthentication no}).
  \item Samba para compartición de ficheros en red \gls{smb}.
  \item Agente Wazuh~4.7.5 con \gls{fim} activo en \texttt{/opt/labprueba}
  con \texttt{realtime="yes"}.
  \item \restic para la gestión de \emph{backups}.
  \item Python~3 con la biblioteca \texttt{cryptography} para el cifrador AES-256-GCM.
  \item \texttt{rsyslog} para que SSH escriba en \texttt{/var/log/auth.log}.
\end{itemize}

El directorio \texttt{/opt/labprueba} contiene 500~ficheros sintéticos generados
con contenido aleatorio de \texttt{/dev/urandom} en cinco tipos (txt, pdf, jpg,
docx, bin) con tamaños entre 1\,KB y 1.100\,KB, produciendo un \emph{dataset} de
aproximadamente 107\,MB no comprimible.

La \textbf{workstation} es la víctima secundaria, objetivo del escenario~E3. Tiene
configuración similar al \emph{fileserver} pero con 300~ficheros en
\texttt{/home/lab/labprueba}.

\subsection{Atacante}

El contenedor \textbf{attacker} se basa en Kali Linux \emph{rolling}. Las
herramientas adicionales instaladas son:

\begin{itemize}
  \item \texttt{awscli}: para las operaciones sobre \gls{s3} en el escenario~E2,
  apuntando al \emph{endpoint} de MinIO con \texttt{--endpoint-url}.
  \item PowerShell~7.4.6 con Atomic Red Team: para la validación de técnicas MITRE
  ATT\&CK (Anexo~B).
  \item \texttt{netcat-openbsd}: para el canal de exfiltración del escenario~E4.
\end{itemize}

El atacante también ejecuta un agente Wazuh, permitiendo monitorizar su actividad
desde el \gls{siem}.

\subsection{Almacenamiento de backups}
\label{sec:almacenamiento-backups}

El contenedor \textbf{restic-sftp} actúa como servidor de destino para la
estrategia~B. Expone únicamente el puerto \gls{ssh} y restringe el acceso al
usuario \texttt{restic-b} mediante autenticación por clave pública.

\textbf{MinIO} implementa un servidor de almacenamiento de objetos compatible con
la \gls{api} de Amazon \gls{s3}, configurado con dos \emph{buckets}:

\begin{description}
  \item[\texttt{restic-c}] Protegido con Object Lock COMPLIANCE durante 30~días.
  El usuario \gls{iam} \texttt{restic-c-user} puede leer y escribir pero no borrar.
  \item[\texttt{restic-b}] Convencional sin protección. El usuario
  \texttt{restic-b-user} tiene permisos estándar de lectura, escritura y borrado.
\end{description}

\section{Diagrama de red y flujo de comunicación}
\label{sec:diagrama-red}

La figura textual~\ref{lst:diagrama-red} resume la topología lógica de la red
\texttt{lab-net} y los flujos de comunicación relevantes para el experimento.
Todos los contenedores comparten la misma red \texttt{bridge} de Docker, lo que
permite resolución de nombres mediante el DNS interno de Docker además de
las IPs estáticas asignadas.

\begin{lstlisting}[style=bash, caption={Topología lógica de la red \texttt{lab-net}
  y flujos de comunicación principales entre contenedores.},
  label={lst:diagrama-red}]
                    +-------------------+
                    |   wazuh-indexer   |  172.20.0.21
                    |   (OpenSearch)    |
                    +---------+---------+
                              ^ 9200 (indexado)
                              |
                    +---------+---------+
                    |   wazuh-manager   |  172.20.0.20
                    | (motor de reglas) |
                    +----^----------^---+
                1514/TCP |          | 1514/TCP  (eventos de agentes)
             (agente FIM)|          |(agente FIM)
        +----------------+          +----------------+
        |                                             |
+-------+-------+                           +---------+-------+
|  fileserver    | 172.20.0.10               |  workstation     | 172.20.0.11
|  (victima E1,E2)|                          |  (victima E3)    |
+---+--------+---+                           +------------------+
    |        |  SSH (E3, lateral)                    ^
    |        +---------------------------------------+
    |  restic backup (B: SFTP / C: S3 API)
    v
+---+-------------+                         +-------------------+
|  restic-sftp     | 172.20.0.41             |      minio        | 172.20.0.40
|  (estrategia B)  |                         | (estrategia C,    |
+------------------+                         |  Object Lock)     |
                                              +--------------------+
        ^ SSH / awscli s3                             ^ awscli s3 (bloqueado
        |  (ataque E2: rm -rf backups)                | por Object Lock)
        +--------------------+-------------------------+
                              |
                    +---------+---------+
                    |     attacker      |  172.20.0.30
                    | (Kali + ART + AWS |
                    |  CLI + netcat)    |
                    +--------------------+
\end{lstlisting}

Se distinguen tres tipos de flujo. El primero es el flujo de \emph{monitorización},
representado por las conexiones \gls{tls} mutuas de los agentes Wazuh del
\emph{fileserver}, la \emph{workstation} y el propio \emph{attacker} hacia el
puerto~1514 del \emph{manager}; este flujo está siempre activo con independencia
del escenario ejecutado. El segundo es el flujo de \emph{backup}, que conecta
periódicamente el \emph{fileserver} con \texttt{restic-sftp} (estrategia~B) o
\texttt{minio} (estrategia~C) durante la fase de creación del \emph{snapshot}.
El tercero es el flujo de \emph{ataque}, iniciado siempre desde el contenedor
\textbf{attacker} hacia el objetivo correspondiente a cada escenario: el
\emph{fileserver} en E1 y E2, la \emph{workstation} en E3 (a través de un salto
SSH intermedio por el \emph{fileserver}), y de nuevo el \emph{fileserver} como
origen de los datos exfiltrados en E4.

Esta separación conceptual en tres flujos es la que permite que las reglas de
detección descritas en la sección~\ref{sec:reglas} del capítulo~\ref{C:disenyo}
puedan asociar de forma no ambigua cada patrón de tráfico o evento \gls{fim} con
su escenario correspondiente, sin necesidad de inspeccionar el contenido de los
paquetes de red.

\section{Recursos, límites y comprobaciones de estado}
\label{sec:recursos-healthchecks}

Cada servicio del fichero \texttt{docker-compose.yml} define explícitamente
límites de memoria (\texttt{mem\_limit}) y una directiva \texttt{healthcheck}
propia, siguiendo la práctica recomendada para entornos de contenedores en los que
la ausencia de límites puede provocar que un único servicio agote la memoria
disponible del \emph{host} y degrade al resto del \emph{stack}. La tabla
resumen de asignación de memoria es la siguiente: 2\,GB para
\texttt{wazuh-indexer} (más 1\,GB de \emph{heap} JVM dentro de ese límite),
1\,GB para \texttt{wazuh-manager}, 512\,MB para \texttt{wazuh-dashboard},
512\,MB para \texttt{minio} y 256\,MB para cada uno de los contenedores
víctima y el atacante.

Las comprobaciones de estado (\texttt{healthcheck}) determinan tanto el orden de
arranque descrito en la sección~\ref{sec:orden-arranque} como la capacidad del
\emph{script} \texttt{verify\_lab.sh} de confirmar automáticamente, antes de
lanzar el experimento, que los nueve contenedores están operativos. Para
\texttt{wazuh-indexer} la comprobación consulta el \emph{endpoint}
\texttt{/\_cluster/health} de OpenSearch; para \texttt{wazuh-manager}, el
\emph{script} de comprobación interno de Wazuh (\texttt{/var/ossec/bin/wazuh-control status});
y para \texttt{minio}, el \emph{endpoint} \texttt{/minio/health/live}. Ningún
contenedor víctima ni el atacante define \texttt{healthcheck} propio, ya que su
disponibilidad se verifica indirectamente a través del estado \emph{Active} del
agente Wazuh correspondiente, reportado por la \gls{api} REST del \emph{manager}.

\section{Mecanismo de comunicación runner.jsonl}
\label{sec:runner}

El volumen Docker \texttt{lab-logs} es el mecanismo de comunicación asíncrona entre
el contenedor atacante y el \emph{manager} Wazuh. Durante la ejecución de cada
escenario, el \emph{script} escribe eventos \gls{json} estructurados en
\texttt{/var/log/lab/\allowbreak{}runner.jsonl}, que el \emph{manager} procesa como
fuente de \emph{logs} mediante \texttt{localfile} con formato \gls{json}.

\begin{lstlisting}[style=bash, caption={Ejemplo de eventos JSON en runner.jsonl
  generados por el escenario E2.}, label={lst:runner-json}]
{"ts":1784804682802708992,"scenario":"E2","event":"phase2_sftp_delete_attempt"}
{"ts":1784804689210081455,"scenario":"E2","event":"immutability_verified",
 "data":"object_lock_blocked_delete"}
\end{lstlisting}

Este mecanismo permite la detección de eventos de alto nivel con independencia de
\texttt{auditd}, con semántica de negocio directamente interpretable y bajo volumen
(5--15 eventos por ejecución).

\section{Configuración de red y seguridad}

Todos los puertos del laboratorio se exponen al \emph{host} únicamente en la
interfaz de \emph{loopback} (\texttt{127.0.0.1}), evitando exposiciones en la red
local.

\begin{table}[ht]
  \centering
  \caption{Puertos expuestos al \textit{host} por el laboratorio.}
  \label{tab:puertos}
  \rowcolors{2}{tabla-par}{white}
  \begin{tabularx}{\textwidth}{C{2.5cm} C{2cm} C{2cm} X}
    \toprule
    \rowcolor{tabla-cabecera}
    \textcolor{white}{\textbf{Servicio}} &
    \textcolor{white}{\textbf{Puerto}} &
    \textcolor{white}{\textbf{Protocolo}} &
    \textcolor{white}{\textbf{Descripción}} \\
    \midrule
    wazuh-manager   & 1514  & TCP   & Recepción de eventos de agentes \\
    wazuh-manager   & 1515  & TCP   & Enrollment de nuevos agentes \\
    wazuh-manager   & 55000 & TCP   & API REST de Wazuh \\
    wazuh-dashboard & 443   & HTTPS & Interfaz web del SIEM \\
    minio           & 9000  & HTTP  & API S3 de MinIO \\
    minio           & 9001  & HTTP  & Consola web de administración \\
    \bottomrule
  \end{tabularx}
\end{table}

\section{Estrategia de registro y trazabilidad}
\label{sec:logging}

La capacidad de depurar el comportamiento del laboratorio y de auditar a
posteriori cada una de las 120~ejecuciones depende de una estrategia de
registro (\emph{logging}) coherente entre los distintos componentes. El
laboratorio combina tres niveles de registro con propósitos distintos y
complementarios.

El primer nivel es el registro de \emph{infraestructura}, generado por Docker
Compose y accesible mediante \texttt{docker compose logs}, que captura la
salida estándar de cada contenedor, incluyendo los mensajes de arranque de
OpenSearch, Wazuh y MinIO. Este nivel es el que se consulta en primer lugar ante
un fallo de arranque del \emph{stack}, como los descritos en la
sección~\ref{sec:limitaciones} siguiente.

El segundo nivel es el registro de \emph{seguridad}, compuesto por los eventos
que los agentes Wazuh envían al \emph{manager} y que este archiva
completamente --- no solo los que generan alerta, gracias a la configuración
\texttt{logall\_json=yes} ya mencionada --- en \texttt{/var/ossec/logs/archives/}
del contenedor \texttt{wazuh-manager}. Este archivo completo, distinto del
subconjunto de eventos que finalmente dispara una regla y llega al
\emph{indexer} como alerta, fue la herramienta principal durante la fase de
calibración de umbrales de frecuencia descrita en la
sección~\ref{sec:reglas} del capítulo~\ref{C:disenyo}, ya que permite
comparar el volumen total de eventos \gls{fim} generados por un escenario
frente al subconjunto que efectivamente supera los umbrales configurados.

El tercer nivel es el registro de \emph{experimento}, correspondiente al
mecanismo \texttt{runner.jsonl} descrito en la
sección~\ref{sec:runner}, que constituye la única fuente de verdad sobre el
progreso interno de cada escenario de ataque (por ejemplo, distinguir si una
ejecución de E3 falló en el paso de acceso SSH inicial o en el salto lateral
posterior). Este tercer nivel es el que consulta directamente
\texttt{experiment\_runner.py} para calcular $\tlaunch$ con precisión de
nanosegundo, como se describe en la sección~\ref{sec:variables-escenario} del
capítulo~\ref{C:disenyo}, y es también la fuente que el mecanismo de control
de errores descrito en la sección~\ref{sec:control-errores} de ese mismo
capítulo utiliza para diferenciar un fallo de infraestructura de un resultado
experimental legítimo.

La combinación de los tres niveles permite reconstruir de forma completa, para
cualquiera de las 120~ejecuciones del experimento, la secuencia exacta de
eventos desde el lanzamiento del escenario hasta el registro final en
\texttt{raw.csv}, lo que resultó fundamental durante la fase de calibración
para diagnosticar por qué ciertas reglas no disparaban en el orden esperado.

\section{Limitaciones conocidas del entorno}
\label{sec:limitaciones}

\subsection{Incompatibilidad de auditd con el kernel CachyOS}

El \emph{kernel} linux-cachyos bloquea el acceso al subsistema \emph{netlink audit}
desde \emph{namespaces} de contenedor incluso con \texttt{privileged:~true}. Esto
impidió el uso de \texttt{auditd} para la instrumentación de \emph{syscalls}. La
solución adoptada fue sustituirlo por \gls{fim} en tiempo real más eventos \gls{json}
estructurados en \texttt{runner.jsonl}, representativa de los despliegues \gls{siem}
en producción donde \texttt{auditd} raramente se activa por su \emph{overhead}.

\subsection{Imposibilidad de webhook de fichero en MinIO}

MinIO no permite escribir eventos de auditoría directamente a un fichero local.
Se resolvió instrumentando los \emph{scripts} de escenario con eventos \gls{json}
en \texttt{runner.jsonl}.

\subsection{Gestión de locks de Restic con Object Lock}

Con Object Lock COMPLIANCE activo, \restic no puede eliminar sus ficheros de
bloqueo. Se usa \texttt{--no-lock} en todas las operaciones con
\texttt{restic-c-user}.

\subsection{rsyslog no disponible por defecto}

Las imágenes mínimas de Ubuntu~24.04 no incluyen \texttt{rsyslog}. Se resolvió
incluyéndolo en los Dockerfiles.

\section{Gestión de certificados TLS y PKI interna}
\label{sec:pki}

La comunicación entre los agentes Wazuh desplegados en los contenedores víctima
y de ataque, y el \emph{manager} central, se protege mediante \gls{tls} mutuo:
tanto el \emph{manager} como cada agente presentan un certificado firmado por
una autoridad de certificación (\gls{ca}) propia del laboratorio, y ambos
extremos verifican el certificado de la contraparte antes de establecer la
conexión sobre el puerto~1514 descrito en la
tabla~\ref{tab:contenedores}. Esta \gls{ca} se genera íntegramente dentro del
propio laboratorio mediante el \emph{script} \texttt{generate\_certs.sh}
(listado~\ref{lst:reproduccion} del anexo~C), sin depender de ninguna entidad
certificadora externa, lo que es coherente con el requisito de reproducibilidad
completa sin credenciales de terceros ya justificado en la
sección~\ref{sec:destruccion-backups} del capítulo~\ref{C:estado_arte} para la
elección de MinIO.

La jerarquía de certificados generada consta de tres niveles: una \gls{ca} raíz
autofirmada, un certificado de servidor para \texttt{wazuh-indexer} y
\texttt{wazuh-manager} firmado por esa \gls{ca}, y certificados de cliente
individuales para cada agente (\emph{fileserver}, \emph{workstation},
\emph{attacker}), también firmados por la misma \gls{ca}. Todos los
certificados se generan con una validez de un año, adecuada para un entorno de
laboratorio de vida corta pero indicativa de que un despliegue en producción de
esta misma arquitectura requeriría un proceso de rotación periódica de
certificados no implementado en el laboratorio actual.

La comunicación entre el \emph{manager} y el \emph{indexer} OpenSearch sigue el
mismo patrón mediante el módulo \texttt{opensearch-security}, cuya
inicialización mediante \texttt{securityadmin.sh} --- detallada en el
listado~\ref{lst:reproduccion} del anexo~C --- es precisamente el paso más
delicado del arranque del laboratorio: debe ejecutarse una única vez, después de
que el \emph{indexer} esté operativo pero antes de que el resto del
\emph{stack} intente conectarse, lo que justifica la división en dos etapas del
arranque descrita en el listado de reproducción del anexo~C y en la
sección~\ref{sec:orden-arranque} siguiente.

\section{Orden de arranque y dependencias}
\label{sec:orden-arranque}

Docker Compose gestiona las dependencias mediante \texttt{depends\_on} con
\texttt{condition:~\allowbreak{}service\_healthy}:

\begin{enumerate}
  \item El \textbf{indexer} arranca y espera hasta que el puerto~9200 responde.
  \item El \textbf{manager} arranca cuando el indexer está \emph{healthy}.
  \item El resto de contenedores arrancan en paralelo cuando el manager está listo.
  \item MinIO y \textbf{minio-init} siguen una cadena de dependencia independiente.
\end{enumerate}

El \emph{script} \texttt{post\_start.sh} realiza la inicialización del
\emph{plugin} de seguridad de OpenSearch, la creación de grupos de agentes y la
verificación de que todos los agentes están en estado \emph{Active}. El tiempo
total de arranque es de aproximadamente tres minutos.

\end{document}
